\subsection{DEFINITION OF THE UNION AND THE INTERSECTION OF A FAMILY OF SETS}

Let X be a family (§3, no. 4) and I its index set. To help the intuitive interpretation of what follows, we shall refer to X as a family of sets. If (X, I, G) is a family of subsets of a set E (that is, a family whose target G is such that the relation Y ∈ G implies Y ⊂ E), we shall denote the family by (X\_t)\_\{t ∈ I\}(X\_t ∈ G) or simply by (X\_t)\_\{t ∈ I\} (§3, no. 6); by abuse of notation we shall denote any family of sets, which has I as its index set, by (X\_t)\_\{t ∈ I\}.

Since the relation \((\forall x)((\iota \in I\) and \(x\in X_{\iota})\Longrightarrow (x\in X_{\iota}))\) is true, it follows from S5 (Chapter I, §4, no 2.) that the relation

\[
(\forall \iota) (\exists Z) (\forall x) ((\iota \in I \text { and } x \in X _ {\iota}) \Longrightarrow (x \in Z))
\]

is true. By virtue of the scheme S8 (§1, no. 6) the relation \((\exists t)(t \in I\) and \(x \in X_t)\) is collectivizing in \(x\) .

DEFINITION 1. Let \((\mathbf{X}_i)_{i\in I}\) be a family of sets (resp. a family of subsets of a set E). The set \(\mathcal{E}_x((\exists i)(i\in I\) and \(x\in \mathbf{X}_i))\) , that is to say, the set of all \(x\) which belong to at least one set of the family \((\mathbf{X}_{\iota})_{\iota \in \mathbf{I}}\) , is called the union of the family, and is denoted by \(\bigcup \mathbf{X}_{\iota}(*)\) .

If \((\mathbf{X}_t)_{i\in I}\) is a family of subsets of a set \(E\) , then its union is a subset of \(E\) ; notice that it does not depend on \(E\) , nor on the target \(\mathfrak{G}\) of the mapping \(t\to X_t\) .

It is clear that if \(\mathbf{I} = \emptyset\) , we have \(\bigcup_{i \in \mathbf{I}} \mathbf{X}_i = \emptyset\) , because the relation \((\exists i)(i \in \mathbf{I}\) and \(x \in \mathbf{X}_i)\) is then false.

Suppose now that \(I \neq \varnothing\) . If \(\alpha\) is an element of I, the relation

\[
(\forall \iota) ((\iota \in I) \Longrightarrow (x \in X _ {\iota}))
\]

implies \(x \in \mathbf{X}_{\alpha}\) and therefore, by virtue of C52 (§1, no. 6), this relation is collectivizing in \(x\) .

DEFINITION 2. Let \((\mathbf{X}_{\iota})_{\iota \in \mathbf{I}}\) be a family of sets whose index set I is not empty. The set \(\mathcal{E}_x((\forall \iota)((\iota \in \mathbf{I}) \Longrightarrow (x \in \mathbf{X}_{\iota}))\) , that is to say, the set of all \(x\) which belong to every set of the family \((\mathbf{X}_{\iota})_{\iota \in \mathbf{I}}\) , is called the intersection of the family and is denoted by \(\bigcap_{\iota} \mathbf{X}_{\iota}\) .

If \(I = \varnothing\) , the relation \((\forall t)((t \in I) \Longrightarrow (x \in X_t))\) is not collectivizing in \(x\) ; for it is a true relation and there exists no set \(Y\) such that \(x \in Y\) is a true relation, because \(Y\) would then be the set of all objects (cf.~§1, no. 7, Remark).

If \((\mathbf{X}_t)_{t \in \mathbf{I}}\) is a family of subsets of a set \(E\) and if \(I \neq \emptyset\) , the relation '' \(x \in E\) and \((\forall t)((t \in I) \Longrightarrow (x \in X_t))\) is equivalent to

\[
(\forall \iota) ((\iota \in I) \Longrightarrow (x \in X _ {\iota}));
\]

consequently it is collectivizing in \(x\) , and the set of all \(x\) which satisfy this relation is equal to \(\bigcap_{i \in I} X_i\) . If \(I = \emptyset\) , the relation `` \(x \in E\) and \((\forall i)((i \in I) \Longrightarrow (x \in X_i))\) '' is equivalent to \(x \in E\) ; it is therefore collectivizing in \(x\) , and the set of all \(x\) which satisfy this relation is \(E\) . Hence we may state the following definition:

DEFINITION 3. Let \((\mathbf{X}_{\iota})_{\iota \in \mathbf{I}}\) be a family of subsets of a set E. The set

\[
\mathcal {E} _ {x} (x \in \mathrm{E} \text {   and   } (\forall \iota) ((\iota \in \mathrm{I}) \Longrightarrow (x \in \mathrm{X} _ {\iota}))),
\]

in other words, the set of all \(x\) which belong to \(\mathbf{E}\) and to each of the sets \(\mathbf{X}_t\) , is called the intersection of the family and is denoted by \(\bigcap_{i=1}^{\infty} \mathbf{X}_t\) .

Hence for a family \((\mathbf{X}_{i})_{i\in\theta}\) of subsets of E we have

\[
\bigcap_ {i \in \theta} X _ {i} = E.
\]

But for a family \((\mathbf{X}_t)_{i\in I}\) of subsets of E whose index set is not empty, the intersection \(\bigcap_{i\in I}\mathbf{X}_t\) depends neither on E nor on the target of \(\iota \to \mathbf{X}_t\) ; and this justifies the use of the same notation in Definitions 2 and 3.

PROPOSITION 1. Let \((\mathbf{X}_{\mathfrak{t}})_{\mathfrak{t}\in \mathbf{I}}\) be a family of sets, and let \(f\) be a mapping of a set K onto I. Then

\[
\bigcup_ {x \in \mathbf {K}} X _ {f (x)} = \bigcup_ {i \in I} X _ {i},
\]

and, if \(\mathbf{I} \neq \emptyset\) ,

\[
\bigcap_ {x \in K} X _ {f (x)} = \bigcap_ {i \in I} X _ {i}.
\]

Let \(x\) be an element of \(\bigcap_{i \in I} X_i\) . There exists an index \(i \in I\) such that \(x \in X_i\) . Since \(f\langle K \rangle = I\) , there exists an index \(x \in K\) such that \(i = f(x)\) , whence \(x \in X_{f(x)}\) and consequently

\[
x \in \bigcup_ {x \in \mathbf {K}} X _ {f (x)}.
\]

Conversely, if \(x \in \bigcup_{x \in \mathbb{K}} X_{f(x)}\) , there exists \(x \in \mathbb{K}\) such that \(x \in X_{f(x)}\) , and therefore, since \(f(x) \in I\) , \(x \in \bigcup_{i' \in I} X_i\) . Hence

\[
\bigcup_ {x \in \mathbb {K}} X _ {f (x)} = \bigcup_ {i \in I} X _ {i}.
\]

Now suppose that \(\mathbf{I} \neq \emptyset\) , and let \(x\) be an element of \(\bigcap_{i \in \mathbf{I}} \mathbf{X}_i\) . For each element \(x\) of \(\mathbf{K}\) we have \(f(x) \in \mathbf{I}\) , hence \(x \in \mathbf{X}_{f(x)}\) , and therefore

\[
x \in \bigcap_ {x \in \mathbb {K}} X _ {f (x)}.
\]

Conversely, let \(x\) be an element of \(\bigcap_{x \in \mathbf{K}} X_{f(x)}\) . If \(\iota\) is any element of \(I\) , there exists an element \(x\) of \(K\) such that \(\iota = f(x)\) , whence \(x \in X_t\) and

consequently \(x \in \bigcap_{i \in I} X_i\) . Hence

\[
\bigcap_ {x \in \mathbf {K}} X _ {f (x)} = \bigcap_ {i \in I} X _ {i}.
\]

For families of subsets of a given set, it is clear that the second part of Proposition 1 remains valid without the restriction \(I \neq \emptyset\) .

COROLLARY. Let \((\mathbf{X}_t)_{i\in \mathbf{I}}\) be a family of sets such that \(\mathbf{X}_t = \mathbf{X}_x\) for each pair of indices \((\mathfrak{t},\mathfrak{x})\) . Then for each \(\alpha \in \mathbf{I}\) we have

\[
\bigcup_ {i \in I} X _ {i} = X _ {\alpha}, \quad \text{and} (\text{if } I \neq \varnothing) \quad \bigcap_ {i \in I} X _ {i} = X _ {\alpha}.
\]

Apply Proposition 1 to the constant mapping \(\iota\to\alpha\) of I onto \(\{\alpha\}\) .

DEFINITION 4. Let \(\mathfrak{F}\) be a set of sets and let \(\Phi\) be the family of sets defined by the identity mapping of \(\mathfrak{F}\) . The union of the sets of \(\Phi\) and (if \(\mathfrak{F}\) is not empty) the intersection of the sets of \(\Phi\) are called, respectively, the union and the intersection of the sets of \(\mathfrak{F}\) , and are denoted by \(\bigcup_{X\sim X^{\prime}}X\) and \(\bigcap_{X\sim X^{\prime}}X\) .

It follows immediately from Proposition 1 that if \((\mathbf{X}_i)_{i\in I}\) is a family of sets, then the union and (if \(I\neq \emptyset\) ) the intersection of the family are respectively equal to the union and the intersection of the sets of the set of elements of this family.

